\documentclass[11pt]{article}
\usepackage[utf8]{inputenc}
\usepackage[T1]{fontenc}
\usepackage{lmodern}
\usepackage[margin=1in]{geometry}
\usepackage{amsmath,amssymb,amsthm}
\usepackage{enumitem}
\usepackage{booktabs}
\usepackage{xcolor}
\usepackage{tikz}
\usepackage[colorlinks=true,bookmarks=false,linkcolor=blue!50!black,urlcolor=blue!50!black]{hyperref}

\theoremstyle{plain}
\newtheorem{theorem}{Theorem}
\newtheorem{lemma}[theorem]{Lemma}
\newtheorem{claim}[theorem]{Claim}
\newtheorem{corollary}[theorem]{Corollary}
\theoremstyle{definition}
\newtheorem{definition}[theorem]{Definition}
\newtheorem{example}[theorem]{Example}
\theoremstyle{remark}
\newtheorem*{remark}{Remark}

\newcommand{\N}{\mathbb{N}}
\newcommand{\Z}{\mathbb{Z}}
\newcommand{\R}{\mathbb{R}}
\newcommand{\Q}{\mathbb{Q}}
\newcommand{\iv}[1]{[#1]}                 % Iverson bracket
\newcommand{\powerset}{\mathcal{P}}
\DeclareMathOperator{\ord}{ord}
\DeclareMathOperator{\gcdop}{gcd}

\setlength{\parindent}{0pt}
\setlength{\parskip}{0.6em}
\newcommand{\lesson}[2]{\href{../lessons/#1}{#2}}   % link to a lesson image

% Let long inline lists like \{1,2,5,6\} break across lines after a comma.
\mathchardef\mycomma="613B
\mathcode`,="8000
{\catcode`,=\active \gdef,{\mycomma\penalty0\hspace{0pt}}}
\emergencystretch=2em

\DeclareMathOperator{\lcmop}{lcm}
\title{Euler, the Basel sum, the Euler product,\\ and coprime probability}
\author{Notes from Carlos' lessons}
\date{}
\begin{document}
\maketitle

\textbf{Source lessons:}
\lesson{2026-06-07-generating-function-euler-sum-of-1-over-n-squared.png}{06-07 (Euler: $\sum1/n^2$)},
\lesson{2026-06-12-generating-function-prob-2-random-chosen-numbers.png}{06-12 (probability two random integers are coprime)}.

\section{The Basel problem: $\displaystyle\sum_{n\ge1}\frac1{n^2}=\frac{\pi^2}6$ (Euler, 1734)}

\paragraph{Euler's idea (board).} Let $f(x)=\dfrac{\sin x}{x}$ (with $f(0):=1$). Its roots are $\pm\pi,\pm2\pi,\pm3\pi,\dots$. Treat $f$ as an ``infinite polynomial'' with those roots and constant term $1$:
\[
\frac{\sin x}x=\Bigl(1-\frac x\pi\Bigr)\Bigl(1+\frac x\pi\Bigr)\Bigl(1-\frac x{2\pi}\Bigr)\Bigl(1+\frac x{2\pi}\Bigr)\cdots=\prod_{n\ge1}\Bigl(1-\frac{x^2}{n^2\pi^2}\Bigr).
\]

\paragraph{Compare the coefficients of $x^2$.}
\begin{itemize}[nosep]
  \item From the Taylor series $\sin x=x-\frac{x^3}{3!}+\frac{x^5}{5!}-\cdots$ (valid for all real $x$): $\frac{\sin x}x=1-\frac{x^2}{3!}+\frac{x^4}{5!}-\cdots$, so $[x^2]\frac{\sin x}x=-\frac16$.
  \item From the product, the $x^2$ coefficient is $-\frac1{\pi^2}\bigl(1+\frac1{2^2}+\frac1{3^2}+\cdots\bigr)$, since only one factor can contribute its $x^2$ term.
\end{itemize}
Equating, $-\frac16=-\frac1{\pi^2}\sum\frac1{n^2}$, so
\[
\boxed{\sum_{n=1}^\infty\frac1{n^2}=\frac{\pi^2}6}.
\]

\paragraph{Caveat.} This is Euler's original \emph{heuristic}: a polynomial is determined by its roots and its value at $0$, but an entire function need not be. It is made rigorous by the Weierstrass/Hadamard factorization, or by other proofs such as the Fourier-series one; the board treats it intuitively, as Euler did.

Zeta notation on the board: $\zeta(s)=\sum_{n\ge1}n^{-s}$. Then $\zeta(1)=+\infty$ (harmonic series); $\zeta(r)<\infty$ for every $r>1$; $\zeta(\tfrac12)=\infty$; $\zeta(2)=\pi^2/6$; $\zeta(2k)$ is a rational multiple of $\pi^{2k}$; and $\zeta(3)=\sum n^{-3}$ is \emph{not} of that form (no closed form in $\pi$).

\section{The Euler product}

\begin{theorem}
For $s>1$, $\displaystyle\zeta(s)=\prod_{p\text{ prime}}\frac1{1-p^{-s}}$.
\end{theorem}
\begin{proof}
For $|y|<1$, $\sum_{k\ge0}y^k=\frac1{1-y}$ (the geometric series; the board's ``For $|x|<1$, $\sum x^k=\frac1{1-x}$''). Multiply these over the primes $p=2,3,5,\dots$ with $y=p^{-s}$:
\[
\Bigl(\sum_{k\ge0}2^{-ks}\Bigr)\Bigl(\sum_{k\ge0}3^{-ks}\Bigr)\Bigl(\sum_{k\ge0}5^{-ks}\Bigr)\cdots=\sum_{(e_2,e_3,e_5,\dots)}\bigl(2^{e_2}3^{e_3}5^{e_5}\cdots\bigr)^{-s}.
\]
By unique prime factorization every $n\ge1$, $n=p_1^{a_1}\cdots p_\ell^{a_\ell}$, appears exactly once as $2^{e_2}3^{e_3}\cdots$. The board expands the first two factors by hand: $\bigl(1+\frac1{2^s}+\frac1{4^s}+\cdots\bigr)\bigl(1+\frac1{3^s}+\frac1{9^s}+\cdots\bigr)$ contains $\frac1{n^s}$ for all $n$ whose only prime factors are $2,3$, and so on. So the product equals $\sum n^{-s}$ (for $s>1$ everything converges absolutely, which justifies the rearrangement).
\end{proof}

At $s=2$: $\zeta(2)=\Bigl(\dfrac1{1-2^{-2}}\Bigr)\Bigl(\dfrac1{1-3^{-2}}\Bigr)\Bigl(\dfrac1{1-5^{-2}}\Bigr)\cdots=\dfrac{\pi^2}6$ (the board's circled result), and $\frac{\pi^2}6\notin\Q$.

\section{The probability that two random integers are coprime is $6/\pi^2$}

The green note on the 06-07 board asks: take two integers $a,b$ ``at random'' (the board admits that ``random'' is not yet defined). The probability that they are coprime is $\frac6{\pi^2}=\frac1{\zeta(2)}$.

\subsection{Heuristic via the Euler product (06-12 board, right side)}
\begin{enumerate}[nosep]
  \item For a random integer $n$, $P[p\mid n]=P[n\equiv0\pmod p]=\frac1p$.
  \item $a,b$ are coprime iff no prime divides both. $P[p\mid a\wedge p\mid b]=\frac1{p^2}$, so $P[p\nmid a\vee p\nmid b]=1-\frac1{p^2}$.
  \item Treating different primes as independent,
  \[
  P[\gcd(a,b)=1]=\prod_p\Bigl(1-\frac1{p^2}\Bigr)=\frac1{\prod_p(1-p^{-2})^{-1}}=\frac1{\zeta(2)}=\frac6{\pi^2},
  \]
  which is the board's boxed $\sum\frac1{n^2}=\frac{1}{(1-\frac1{2^2})(1-\frac1{3^2})\cdots}$ (read through the reciprocal) combined with Section~1.
\end{enumerate}
The independence step is informal: there is no uniform probability measure on $\N$, so the claim is rigorous only as a statement about limiting frequencies in $[1,n]^2$ (see the next subsection and the references).

\subsection{Counting via Euler's totient (06-12 board, left side)}

Let $P(n)=\dfrac{\bigl|\{(a,b)\in[n]^2:\gcd(a,b)=1\}\bigr|}{n^2}$. We want $\lim P(n)$.

\paragraph{Pair count.} Write $\varphi(m)=|\{k\in[m]:\gcd(k,m)=1\}|=\sum_{k=1}^m\iv{\gcd(k,m)=1}$. Pairs with $a\le b$ and $\gcd=1$: for each $b\ge2$ there are $\varphi(b)$ choices of $a$ (and $a=b$ is excluded since $\gcd(b,b)=b>1$); for $b=1$ there is the single pair $(1,1)$. By the symmetry $a\leftrightarrow b$ the total is
\[
1+2\sum_{b=2}^n\varphi(b)=2\sum_{b=1}^n\varphi(b)-1\qquad(\text{since }\varphi(1)=1).
\]
The board shows exactly this ($\frac1{n^2}\bigl(1+2\sum_{b=2}^n\varphi(b)\bigr)$). Check for $n=5$: $\varphi=1,1,2,2,4$ has sum $10$, giving $19$, which matches a direct count of the coprime pairs in $[5]^2$.

\paragraph{Facts about $\varphi$ (board).}
\begin{itemize}[nosep]
  \item $\varphi(p^k)=p^k-p^{k-1}=p^{k-1}(p-1)$ (remove the multiples of $p$).
  \item $\varphi(ab)=\varphi(a)\varphi(b)$ if $\gcd(a,b)=1$ (Chinese remainder theorem).
  \item Hence for $n=p_1^{a_1}\cdots p_\ell^{a_\ell}$: $\varphi(n)=n\prod_i\bigl(1-\frac1{p_i}\bigr)$.
  \item $\sum_{d\mid n}\varphi(d)=n$.
\end{itemize}
\begin{proof}[Proof of the last fact]
Sort $\{1,\dots,n\}$ by $g=\gcd(k,n)$. Those $k$ with $\gcd(k,n)=n/d$ correspond to $k/(n/d)\in[d]$ coprime to $d$, so there are $\varphi(d)$ of them. Summing over $d\mid n$ counts all $n$ numbers.
\end{proof}

\paragraph{The limit} (my completion; the board stops after setting up $S(n)=\sum_{k\le n}\varphi(k)$). M\"obius inversion of $\sum_{d\mid n}\varphi(d)=n$ gives $\varphi(k)=\sum_{d\mid k}\mu(d)\,\frac kd$, so
\[
S(n)=\sum_{k\le n}\varphi(k)=\sum_{d\le n}\mu(d)\sum_{m\le n/d}m
=\sum_{d\le n}\mu(d)\,\frac{\lfloor n/d\rfloor(\lfloor n/d\rfloor+1)}2
=\frac{n^2}2\sum_{d\ge1}\frac{\mu(d)}{d^2}+O(n\log n),
\]
and $\sum_d\mu(d)d^{-2}=\prod_p(1-p^{-2})=\frac1{\zeta(2)}$ (the Euler product again). So $S(n)\sim\frac{3}{\pi^2}n^2$ and $P(n)=\frac{2S(n)-1}{n^2}\to\frac6{\pi^2}$. Numerically, $S(2000)/2000^2=0.3041$ against $3/\pi^2=0.3040$.

\section*{Asides on these boards}
\begin{itemize}[nosep]
  \item $1+2+3+4+\cdots=-\frac1{12}$ appears (flagged) on the 06-12 board, with $\zeta(k)=S(\infty,-k)=\lim_{n\to\infty}S(n,-k)$ and $S(k,n)=\sum_{i\le n}i^k$: the regularized value $\zeta(-1)=-\frac1{12}$, not a convergent sum.
  \item The 06-12 board also lists derivative rules for formal series: $(A+B)'=A'+B'$, $(AB)'=A'B+AB'$, $(x_0,x_1,x_2,\dots)'=(x_1,2x_2,3x_3,\dots)$; and $\sum nx^n=\frac x{(1-x)^2}$, i.e.\ $(0,1,2,3,\dots)$; and the recurrence $a_{n+1}=2a_n+n$, which is solved in the topic \emph{Ordinary generating functions}.
\end{itemize}

\section*{Sources}
\begin{itemize}[nosep]
  \item \href{https://arxiv.org/pdf/1612.03700}{The probability that two random integers are coprime --- Bureaux \& Enriquez} and \href{https://arxiv.org/pdf/1806.00053}{On the probability that two random integers are coprime --- Lei \& Kadane}: the precise meaning of ``random integer'' (no countably additive uniform measure on $\N$).
  \item \href{https://www.cut-the-knot.org/m/Probability/TwoCoprime.shtml}{Probability of Two Integers Being Coprime (cut-the-knot)}; \href{https://aperiodical.com/2013/06/cushing-your-luck-properties-of-randomly-chosen-numbers/}{Cushing your luck (The Aperiodical)}: $1/\zeta(2)$ and the Basel value.
  \item The totient/M\"obius completion above is standard (e.g.\ Hardy \& Wright, Thm.\ 331, as cited by Lei--Kadane); the numeric checks were run locally.
\end{itemize}

\end{document}
